% ==========================================================================
% CHAPTER 1 -- INTRODUCTION
% ==========================================================================

\chapter{Introduction}
\label{ch:introduction}

Cardiovascular diseases (CVDs) continue to be a leading cause of death
around the world. Among them, cardiac arrhythmias, irregularities of the
heart rate caused by abnormal electrical signals in the heart, can lead to
serious complications such as stroke, heart failure, and sudden cardiac
arrest if they are not detected in time. This thesis develops a deep
learning based computer-aided system for automated ECG arrhythmia
classification and evaluates it under a clinically meaningful,
record-disjoint protocol that is suitable for wearable, resource-constrained
deployment.


\section{Overview}
\label{sec:overview}

The electrocardiogram (ECG) is the most widely used non-invasive diagnostic
tool for assessing cardiac electrophysiology and diagnosing arrhythmias.
However, ECG reading is time consuming and is normally performed manually,
which makes it vulnerable to inter-observer variability and limits its use
for long-term or continuous monitoring. A single 24-hour ambulatory record
can contain on the order of $10^{5}$ heart beats, and annotating such a
record manually is expensive and impractical at scale. There is therefore a
clear and growing need for automated, high-throughput screening systems that
can classify beats rapidly, precisely, and reliably.

Recent advances in deep learning have substantially improved automated ECG
analysis. Convolutional neural networks (CNNs) extract ECG morphology
directly from the raw signal, while recurrent networks such as long
short-term memory (LSTM) architectures capture the temporal relationships
between successive cardiac beats. Self-attention and Transformer encoders
have also been introduced to focus the model on the most informative signal
regions, increasing both classification performance and interpretability.
Channel-recalibration mechanisms such as squeeze-and-excitation (SE) blocks
and feature-wise conditioning schemes like FiLM further improve the ability
of a network to discriminate morphologically similar classes at modest
parameter cost. Finally, explainable artificial intelligence (XAI)
techniques such as Grad-CAM and Integrated Gradients provide visual
explanations of model decisions, which helps build clinician confidence in
automated systems.

In this thesis, a hybrid CNN--BiLSTM--attention model and a lightweight
multi-scale feature-augmented network (MSFT-Net) are designed, trained, and
compared for five-class ECG beat classification under a strict
record-disjoint inter-patient protocol, and the best model is quantised and
deployed on an STM32 microcontroller.


\section{Background and Motivation}
\label{sec:background}

Despite the progress described above, several key issues remain unresolved
in the automated ECG classification literature. Many existing studies focus
on maximising classification accuracy while providing little transparency
into the decision-making process of the deep learning model. Public ECG
datasets are skewed, noisy, and dominated by normal beats, which reduces the
capability of models trained on them to recognise rare but clinically
important abnormal classes.

An even more serious, and less frequently discussed, issue concerns the
evaluation protocol. A large portion of the literature reports results under
a \emph{beat-level} split in which beats from the same patient appear in
both training and test sets. Because consecutive beats of one patient are
strongly correlated in morphology and rhythm, a sufficiently expressive
model can simply memorise the patient-specific beat templates and recall
them at test time. The resulting figures measure within-patient recall
rather than generalisation to unseen patients, and they are not comparable
to inter-patient results, which typically reach much lower performance
levels, especially for minority classes. The ANSI/AAMI EC57 standard
recommends the record-disjoint DS1/DS2 partition of de~Chazal
et~al.~\cite{dechazal2004} precisely to avoid this problem, yet recent
systematic reviews of the field show that both the standard partition and
proper class grouping are frequently bypassed~\cite{silva2025systematic,xiao2023systematic}.

Finally, the five AAMI superclasses are not equally difficult, and each
resistant class fails for different reasons. A supraventricular ectopic beat
($S$) is defined by \emph{prematurity}, a near-normal QRS complex arriving
too early, so a fixed window centred on the R-peak alone cannot separate
inter-patient $S$ beats from normal ($N$) beats without timing information.
A fusion beat ($F$) is, by definition, a morphological mixture of a normal
and a ventricular depolarisation, so its features are inherently ambiguous.
The unclassifiable class ($Q$) consists almost entirely of beats from four
paced records pointed to by EC57, so its statistics rest on very few
patients. These class-specific challenges motivate the feature engineering
and targeted treatments developed in this work.

The motivation of this thesis is therefore threefold: (i) to build a hybrid
deep learning classifier that learns ECG morphology and rhythm together, so
that timing-sensitive classes such as $S$ are recognised correctly;
(ii) to evaluate it honestly under a record-disjoint inter-patient protocol
so that the reported numbers reflect genuine generalisation to unseen
patients; and (iii) to make the model compact and quantisable so that it can
run in real time on a microcontroller-class wearable device, which is the
computing environment available in real ambulatory monitoring.


\section{Literature Review}
\label{sec:literature}

Automated ECG arrhythmia classification research has developed along several
largely independent dimensions (architecture design, class-imbalance
mitigation, edge deployability, explainability, and evaluation
methodology) that are rarely optimised together within a single study. This
review considers each axis in turn, focusing on studies published mainly
between 2020 and 2026 together with the foundational earlier works that
established the field's core methods, and closes by identifying the gap
that motivates the present work.

\subsection{CNN--recurrent hybrids and channel attention}

Convolutional and recurrent architectures remain the predominant building
blocks for ECG beat classification. Kiranyaz~et~al.~\cite{kiranyaz2015}
made one of the first contributions, showing that patient-specific 1-D
convolutional classifiers can run in real time on wearable devices; the
method is not scalable to population level, however, because the filters are
tuned per person. Channel-recalibration mechanisms were later transplanted
into ECG backbones: Zhang~et~al.~\cite{zhang2020seecgnet} proposed
SE-ECGNet, a multi-scale residual network whose SE gating ``reweighs''
convolutional channels and improves discrimination between morphologically
similar beat types at negligible parameter cost. Zheng~et~al.
\cite{zheng2025multiinput} extended this idea with a two-branch, multi-input
CNN that fuses short-window and long-window spectrograms of the same beat,
reporting 99.13\% and 95.84\% accuracy, and 94.46\% and 95.91\% macro-$F_1$,
on the MIT-BIH and SPH databases respectively. A recurring finding of this line of work is
that narrowly directing attention at the channel level does not by itself
close the performance gap between the majority normal class and the
clinically important minority ectopic classes.

\subsection{Residual, spatial, and self-attention architectures}

A second stream of research places attention at the spatial level.
Xu~et~al.~\cite{xu2022racnn} fused residual skip connections with a U-Net
like attention gate (RA-CNN) and reported 98.5\% accuracy on a strict
inter-patient DS1/DS2 split, with class-wise $F_1$ scores of 82.8\% and
91.7\% for the supraventricular and ventricular classes, figures that are
informative precisely because they are reported under the more difficult
inter-patient regime. Full Transformer self-attention has also been applied:
Islam~et~al.~\cite{islam2024} proposed CAT-Net, a cascaded convolutional,
attention, and Transformer network for single-lead ECG classification that
reached 99.14\% accuracy on MIT-BIH by learning long-range inter-beat
dependencies that fixed-receptive-field CNNs cannot represent; the authors
note, however, that the model is too large for wearable deployment. Xia
et~al.~\cite{xia2023transformer} reported a lighter design combining a CNN
front end with a denoising autoencoder and a Transformer encoder, evaluated
explicitly under the inter-patient paradigm, and Anjum~et~al.
\cite{anjum2023temporal} proposed a temporal Transformer-based fusion
framework for morphological (rather than rhythmic) arrhythmia classes.
Taken together, these studies suggest that self-attention is not a
substitute for, but rather a complement to, convolutional morphology
extraction, a finding that directly motivates the dual-branch
CNN--Transformer design considered in this thesis.

\subsection{Hybrid CNN--BiLSTM--attention frameworks and gated fusion}

Recurrent sequence modelling is frequently combined with convolutional and
attention modules to capture morphology and rhythm simultaneously. Mechichi
and Benzarti~\cite{mechichi2025enhanced} proposed a CNN--CBAM--BiLSTM
nested network in which convolutional block attention rectifies spatial and
channel features and a BiLSTM layer models beat-to-beat dependencies,
reporting 99.20\% accuracy, 97.50\% sensitivity, and 98.29\% mean $F_1$ on
MIT-BIH after SMOTE-based rebalancing; the authors explicitly identify the
BiLSTM stage as the main computational burden, a limitation that recurs
throughout the lightweight-deployment literature. The family of
architectures also disagrees on how parallel branches should be integrated:
concatenation, summation, and learned gating are all used, and the choice of
fusion mechanism carries measurable penalties for minority-class recall
rather than being a neutral implementation detail. Feature-wise linear
modulation (FiLM)~\cite{perez2018film} offers a lightweight, per-channel
conditioning alternative for injecting auxiliary information.

\subsection{Image-based and time-frequency representations}

For the three-class separation of arrhythmia (ARR), congestive heart
failure (CHF), and normal sinus rhythm (NSR), converting the 1-D signal into
a 2-D time-frequency image has been fruitful. Olanrewaju~et~al.
\cite{olanrewaju2021cwt} classified scalograms obtained with the continuous
wavelet transform (CWT) using an AlexNet-based deep network, reaching 98.7\%
accuracy on the ARR/CHF/NSR corpus. Madan~et~al.~\cite{madan2022} fused a 2-D CNN
with an LSTM stage running over scalogram sequences to reach similar
accuracy, while Prusty~et~al.~\cite{prusty2024} paired SIFT features with a
CNN classifier to achieve 99.78\% accuracy at the cost of an extra,
non-differentiable feature-extraction stage that complicates end-to-end
optimisation. These studies confirm the discriminative power of image-based
representations for ARR/CHF/NSR, but none addresses microcontroller-class
deployment or a record-disjoint inter-patient partition, a limitation that
is revisited in the evaluation-critique part of this review.

\subsection{Class imbalance and synthetic minority augmentation}

The fusion ($F$) and supraventricular ($S$) classes together contain fewer
than 5\% of the beats in MIT-BIH, so class-imbalance treatment is
near-ubiquitous in the literature. The most popular baseline is SMOTE-family
oversampling~\cite{chawla2002smote}, used directly in several of the works
cited above; more recent generative approaches train conditional or
Wasserstein GANs to synthesise physiologically plausible minority-class
beats, achieving up to 4--9 percentage points of improvement in minority
precision when the synthetic beats are quality-screened against real-beat
templates. A major drawback of these methods is that unscreened synthetic
beats can be trivially distinguished from real ones by a downstream
discriminator; an almost-perfect discriminator AUC is reported as direct
proof of an under-converged generator rather than of class-balancing
success. This diagnostic (training a real-versus-synthetic discriminator
and using its AUC as a quality-control gate) is adopted as a validation
step in the augmentation pipeline of the present study.

\subsection{Lightweight and edge-deployable architectures}

A growing body of literature considers deployment directly on
microcontroller-class hardware rather than treating quantisation as an
afterthought. Busia~et~al.~\cite{busia2024tinytransformer} designed a tiny
Transformer with only 6k parameters that achieves 98.97\% accuracy on five
AAMI classes with sub-8-bit integer inference, executing one inference in
4.28\,ms at 0.09\,mJ on the ultra-low-power GAP9 processor, evidence that
self-attention, previously thought to require too much memory for
microcontrollers, can be made competitive with convolutional baselines. At
the opposite end of the hardware spectrum, Zambrano-de~la~Torre
et~al.~\cite{zambrano2026tinyml} implemented a quantised 1-D CNN on an 8-bit
Arduino UNO (32\,KB flash, 2\,KB SRAM), achieving 97.6\% accuracy on a
reduced three-class problem with a footprint below 24\,KB INT8. The authors
acknowledge, however, that their evaluations use a beat-level split without
patient separation, a known source of overly optimistic numbers. The
three-way trade-off between representational capacity, efficiency, and
feasible evaluation is the central concern of this literature, and it is
what motivates the ``dual-model'' design (an accuracy-oriented hybrid and a
deployment-oriented lightweight network) used in the present study.

\subsection{Evaluation methodology: the inter-patient/intra-patient divide}

Almost all of the studies above share a methodological thread that deserves
separate treatment. Two recently published systematic reviews of the ECG
arrhythmia classification literature, covering the years 2017--2024, found
that a substantial fraction of published works either fail to state whether
their split is inter-patient or intra-patient, or report only the more
favourable intra-patient figure~\cite{silva2025systematic,xiao2023systematic}.
Silva~et~al.~\cite{silva2025systematic} formalise an ``Embedded, Clinical,
and Comparative Criteria'' (E3C) checklist and find that very few surveyed
studies satisfy inter-patient partitioning, AAMI-compliant class grouping,
and embedded feasibility simultaneously. This reporting inconsistency has
real consequences: the 85--100\% accuracy figures reported under beat-level
splits collapse dramatically under a record-disjoint protocol (equivalent
architectures have been shown to drop to below 60\% inter-patient
accuracy) because beats from the same patient are morphologically
correlated and the model exploits this correlation to memorise
patient-specific templates rather than to generalise. The headline accuracy
values reported throughout the literature are therefore not directly
comparable unless the evaluation protocol is controlled for.

\subsection{Explainability}

Clinical adoption of automated ECG classifiers depends on the decision
process being transparent rather than a black box. Grad-CAM, originally
designed for 2-D image classifiers, has been adapted to 1-D ECG backbones to
produce beat-level saliency maps that can be checked against known
P--QRS--T morphology~\cite{selvaraju2017gradcam}. Integrated Gradients
offers a complementary, axiomatically based attribution that does not rely
on the selection of a specific convolutional layer
\cite{sundararajan2017integrated}. Grad-CAM remains the most common method
in the XAI-for-ECG literature and is normally applied to verify
qualitatively whether the classifier's attention falls on the QRS complex
and T-wave rather than on dataset-specific artefacts. Quantitative
validation of this alignment (for example by deletion curves or
intersection-over-union against cardiologist-annotated intervals) is
reported far less often, and this thesis treats it as a faithfulness
challenge rather than assuming correctness from visual inspection alone.

\subsection{Summary and identified gap}

The surveyed literature documents separately that: (i) channel and spatial
attention improve discrimination of morphologically similar
beats~\cite{hu2018squeeze,zhang2020seecgnet,xu2022racnn,zheng2025multiinput};
(ii) long-range dependencies that convolutional or recurrent layers alone
cannot capture are effectively handled by transformer
self-attention~\cite{islam2024,xia2023transformer,anjum2023temporal,vaswani2017attention};
(iii) gated, learnable fusion of parallel branches outperforms naive
concatenation; (iv) quality-screened generative or SMOTE-based augmentation
recovers minority-class recall~\cite{chawla2002smote}; (v) models with
sub-10k parameters can achieve competitive accuracy with INT8 quantisation
on microcontroller-class hardware~\cite{busia2024tinytransformer,zambrano2026tinyml};
and (vi) none of the accuracy figures above is trustworthy unless it is
reported under a protocol that withstands inter-patient
scrutiny~\cite{silva2025systematic,xiao2023systematic,dechazal2004}. What no
single study has combined is a clearly stated record-disjoint protocol,
explicit and ablated treatment of both of the hardest AAMI classes (fusion
and supraventricular), a faithfulness-tested explainability layer, and a
measured, rather than asserted, microcontroller deployment of the full
inter-patient model. Two architectural families (a CNN--BiLSTM--attention
hybrid and a gated-fusion lightweight CNN) are therefore evaluated in this
thesis under identical data splits so that they can be compared
statistically rather than across incommensurable papers. This gap is the
direct driver of the research described in the remainder of this book.


\section{Research Gap}
\label{sec:gap}

Existing automated ECG classification systems suffer from three
interlocking weaknesses. First, most published results are obtained under
beat-level or otherwise leaky splits that inflate accuracy by letting the
model memorise patient-specific templates, so the reported numbers do not
reflect how the system would behave on an unseen patient. Second, the two
most clinically dangerous minority classes (fusion and supraventricular
ectopic beats) remain poorly recognised because they require timing and
template information that raw-waveform-only models do not receive. Third,
most accurate models are too large to run on the microcontroller-class
hardware found in wearable monitors, and the few lightweight models that do
run on such hardware are typically validated under optimistic, non-standard
protocols. Therefore, there is a need for a compact, feature-augmented deep
learning classifier that is trained and tested under a strictly
record-disjoint inter-patient protocol, treats the resistant $S$ and $F$
classes explicitly, verifies its explainability quantitatively, and is
demonstrated to run in real time on an actual microcontroller.


\section{Objective}
\label{sec:objectives}

The objectives of this thesis are:

\begin{enumerate}
  \item To design a hybrid CNN--BiLSTM--attention model and a lightweight
        multi-scale feature-augmented network (MSFT-Net) that classify the
        five ANSI/AAMI EC57 beat classes (N, S, V, F, Q) from one-second
        R-peak-centred ECG windows combined with rhythm, morphological, and
        template-similarity auxiliary features.
  \item To evaluate both architectures under a strictly record-disjoint,
        inter-patient DS1/DS2 protocol (and a permissive intra-patient
        six-category ECG protocol that adds a CHF-derived clinical-condition
        category for comparison), reporting per-class sensitivity
        and positive predictivity with statistical significance testing.
  \item To analyse the contribution of each architectural and feature
        component through a controlled ablation study, and to verify the
        faithfulness of the Grad-CAM and Integrated-Gradients explanations
        using deletion curves.
  \item To quantise the best inter-patient model to INT8 and deploy it on an
        STM32 microcontroller, measuring accuracy retention, model size,
        latency, and memory footprint.
\end{enumerate}


\section{Thesis/Project Outline}
\label{sec:outline}

The remainder of this book is organised as follows.

\textbf{Chapter~\ref{ch:methodology}} describes the methodology adopted in
this work: the MIT-BIH and auxiliary databases, beat extraction and
filtering, the record-disjoint partitioning scheme, the auxiliary feature
vector, the two network architectures, the training and imbalance-handling
procedures, and the statistical evaluation and explainability methods.

\textbf{Chapter~\ref{ch:results}} presents the results obtained: the
exploratory intra-patient six-category ECG study, the inter-patient five-class study with its
AAMI EC57 reporting, the ablation and fusion-beat analyses, the
explainability and robustness results, and the microcontroller deployment
measurements.

\textbf{Chapter~\ref{ch:social}} examines the social and environmental
influence of the work, including its financial, health, safety, legal, and
sustainability implications for Bangladesh.

\textbf{Chapter~\ref{ch:po}} explains how the work addresses the program
outcomes, complex engineering problems and complex engineering activities
defined in the EEE 4000 manual.

\textbf{Chapter~\ref{ch:conclusion}} concludes the book and outlines the
future plan.

The work was carried out over the fourth year of the B.Sc.\ in
Electrical \& Electronic Engineering programme at RUET.